\documentclass[11pt,a4paper]{article}
\usepackage[utf8]{inputenc}
\usepackage[T1]{fontenc}
\usepackage{amsfonts}
\usepackage[french]{babel}
\frenchbsetup{StandardLists=true}
\usepackage{enumitem}
\usepackage{amssymb}
\usepackage{amsmath}
\usepackage[left=2cm,right=2cm,top=2cm,bottom=2cm]{geometry}
\usepackage{multirow}
\usepackage[dvipsnames]{xcolor}

\usepackage[T1]{fontenc}
\usepackage{fouriernc}
\usepackage{graphicx}

\usepackage{setspace}
\singlespacing

\usepackage{tcolorbox}
\usepackage{multicol}
\usepackage{caption}
\usepackage{fancyhdr}
\pagestyle{fancy}
\renewcommand\headrulewidth{1pt}
\fancyhead[L]{Lycée Denis Diderot }
\fancyhead[R]{Classe de 3PM}
\setcounter{page}{1}\fancyfoot[C]{page \thepage}
\fancyfoot[R]{Année 2025-2026}
\fancyfoot[L]{Alexis RUIZ}
\usepackage{multirow,tabularx,booktabs}
\newcommand{\cadreblanc}[1]{%
\medskip\framebox[\linewidth][c]{%
\rule{0mm}{#1mm}}\par
\medskip}

\usepackage{awesomebox}
\setlength{\aweboxvskip}{0mm}
\usepackage{pifont}
\usepackage{chemmacros}
% \chemsetup{modules=all} % obsolète en v6, les modules sont chargés automatiquement
\setlength{\columnseprule}{.25mm}
\setlength{\parindent}{0pt}

\usepackage{tikz}
\usepackage{circuitikz}
\usepackage{pgfplots}
\pgfplotsset{compat=1.18}

\renewcommand{\arraystretch}{2}

\frenchbsetup{StandardLists=true}

\begin{document}

\begin{center}
 \LARGE{Chapitre -- Mesure de l'intensité électrique}
\end{center}

\normalsize

\hrule\
\\[0,2cm]

L'intensité électrique est une grandeur fondamentale en électricité. On la retrouve sur toutes les fiches techniques d'appareils électriques : ampoules, chargeurs, moteurs... Elle permet de caractériser le flux de charges qui circule dans un circuit. Sa mesure correcte est indispensable pour comprendre et dimensionner tout montage électrique.

\vspace{3mm}

\textbf{1. L'intensité électrique} \\

\textbf{1.1. Définition}

\begin{tcolorbox}[colback=red!5!white,colframe=red!75!black]
{
L'\textbf{intensité électrique} est une grandeur physique qui caractérise le \textbf{débit de charges électriques} (d'électrons) circulant dans un circuit.

\begin{itemize}
    \item L'intensité se note \textbf{I}.
    \item Elle se mesure en \textbf{ampère}, de symbole \textbf{A}.
    \item Plus l'intensité est grande, plus les charges circulent rapidement dans le conducteur.
\end{itemize}
}
\end{tcolorbox}

\textbf{Remarque : } L'ampère porte le nom du physicien français \textbf{André-Marie Ampère} (1775--1836), l'un des fondateurs de l'électromagnétisme. \\

\textbf{1.2. Les multiples et sous-multiples}

\vspace{8mm}

\begin{center}
\renewcommand{\arraystretch}{1.5}
\begin{tabular}{|l|c|l|}
\hline
\textbf{Unité} & \textbf{Symbole} & \textbf{Valeur en ampère} \\
\hline
Ampère       & A    & $1\,\text{A}$ \\
\hline
Milliampère  & mA   & $1\,\text{mA} = 10^{-3}\,\text{A} = 0{,}001\,\text{A}$ \\
\hline
Microampère  & $\mu$A & $1\,\mu\text{A} = 10^{-6}\,\text{A}$ \\
\hline
\end{tabular}
\end{center}

\vfill

\textbf{2. L'ampèremètre} \\

\textbf{2.1. Rôle et symbole}

\begin{tcolorbox}[colback=red!5!white,colframe=red!75!black]
{
L'\textbf{ampèremètre} est l'instrument de mesure utilisé pour mesurer l'intensité électrique dans un circuit. Son symbole dans un schéma électrique est un \textbf{cercle contenant la lettre A}.

\begin{center}
\begin{tikzpicture}[scale=1]
  \draw[thick] (0,0) -- (0.35,0);
  \draw[thick] (0.7,0) circle (0.35);
  \node at (0.7,0) {\textbf{A}};
  \draw[thick] (1.05,0) -- (1.75,0);
  \node at (0.875,-0.7) {\small Symbole de l'ampèremètre};
\end{tikzpicture}
\end{center}
}
\end{tcolorbox}

\vfill 

\textbf{2.2. Les bornes de l'ampèremètre}

Un ampèremètre possède généralement \textbf{trois bornes} :

\begin{itemize}
    \item La borne \textbf{COM} (ou \textbf{--}) : borne commune, toujours utilisée, côté négatif du circuit.
    \item La borne \textbf{A} : pour mesurer des intensités jusqu'à \textbf{10 A}.
    \item La borne \textbf{mA} : pour mesurer des intensités inférieures à \textbf{200 mA}, avec plus de précision.
\end{itemize}

\vfill 

\newpage

\warningbox{Ne jamais brancher l'ampèremètre directement aux bornes d'un générateur sans résistance. Cela provoquerait un \textbf{court-circuit} pouvant détruire l'appareil ou endommager le générateur.}

\vfill

\textbf{3. Branchement de l'ampèremètre} \\

\textbf{3.1. Règle de branchement}

\begin{tcolorbox}[colback=red!5!white,colframe=red!75!black]
{
L'ampèremètre se branche \textbf{en série} dans le circuit, c'est-à-dire qu'il est \textbf{intercalé} dans le circuit : le courant électrique doit le traverser. Il ne faut jamais le brancher en parallèle.
}
\end{tcolorbox}

\begin{center}
\begin{tikzpicture}[scale=1.3]
  % === Fil gauche : générateur ===
  \draw[thick] (0,0) -- (0,0.9);
  \draw[thick] (0,1.25) circle (0.35);
  \node at (0,1.25) {\textbf{G}};
  \node[left=5pt] at (-0.35,1.25) {\small Gén.};
  \draw[thick] (0,1.6) -- (0,2.5);
  % === Fil haut ===
  \draw[thick] (0,2.5) -- (3,2.5);
  % === Fil droit : lampe ===
  \draw[thick] (3,2.5) -- (3,1.6);
  \draw[thick] (3,1.25) circle (0.35);
  \draw[thick] (2.75,1.0) -- (3.25,1.5);
  \draw[thick] (2.75,1.5) -- (3.25,1.0);
  \node[right=5pt] at (3.35,1.25) {\small Lampe};
  \draw[thick] (3,0.9) -- (3,0);
  % === Fil bas : ampèremètre ===
  \draw[thick] (0,0) -- (1.0,0);
  \draw[thick] (1.35,0) circle (0.35);
  \node at (1.35,0) {\textbf{A}};
  \node[below=4pt] at (1.35,-0.45) {\small Ampèremètre};
  \draw[thick] (1.7,0) -- (3,0);
  % === Flèche courant ===
  \draw[->, thick, red] (0.3,2.7) -- (1.3,2.7);
  \node[above, red, font=\small] at (0.8,2.7) {$\mathrm{I}$};
\end{tikzpicture}
\end{center}

\begin{center}
\small\textit{Schéma : ampèremètre branché en série avec un générateur et une lampe}
\end{center}

\vfill 

\textbf{3.2. Méthode de branchement}

\begin{tcolorbox}[colback=red!5!white,colframe=red!75!black]
{
\begin{enumerate}
    \item \textbf{Ouvrir} le circuit à l'endroit où l'on souhaite mesurer l'intensité.
    \item Choisir d'abord le \textbf{plus grand calibre} (borne A) pour protéger l'appareil.
    \item Brancher la borne \textbf{COM} du côté \textbf{négatif} du générateur (potentiel bas).
    \item Brancher la borne \textbf{A} du côté \textbf{positif} du générateur (potentiel haut).
    \item Si la valeur lue est trop faible, basculer sur le calibre \textbf{mA} pour plus de précision.
\end{enumerate}
}
\end{tcolorbox}

\vfill

\textbf{4. Lecture de la mesure} \\

\textbf{Choisir le bon calibre}

\begin{tcolorbox}[colback=red!5!white,colframe=red!75!black]
{
Le \textbf{calibre} est la valeur maximale que peut afficher l'ampèremètre sur une gamme donnée. Il faut toujours choisir le \textbf{calibre immédiatement supérieur} à la valeur attendue.
}
\end{tcolorbox}

\textbf{Attention : } Commencer toujours par le plus grand calibre, puis réduire si nécessaire. Si l'écran affiche une erreur ou dépasse la plage, le calibre est trop petit. \\

\vfill

\newpage

\textbf{5. Intensité dans un circuit} \\

\textbf{5.1. Circuit en série}

\begin{tcolorbox}[colback=red!5!white,colframe=red!75!black]
{
Dans un circuit \textbf{en série}, l'intensité est \textbf{identique} en tout point du circuit :
$$\mathrm{I}_1 = \mathrm{I}_2 = \mathrm{I}_3 = \ldots$$
L'ampèremètre peut donc être placé n'importe où dans la boucle.
}
\end{tcolorbox}

\vfill

\textbf{5.2. Circuit en dérivation (loi des nœuds)}

\begin{tcolorbox}[colback=red!5!white,colframe=red!75!black]
{
Dans un circuit \textbf{en dérivation}, l'intensité dans la branche principale est égale à la \textbf{somme} des intensités dans chaque branche dérivée. C'est la \textbf{loi des nœuds} :
$$\mathrm{I} = \mathrm{I}_1 + \mathrm{I}_2$$
\begin{center}
\begin{tikzpicture}[scale=1.1]
  % === Fil haut ===
  \draw[thick] (0,3) -- (4,3);
  % === Fil bas ===
  \draw[thick] (0,0) -- (4,0);
  % === Fil gauche : générateur disque G ===
  \draw[thick] (0,0) -- (0,1.1);
  \draw[thick] (0,1.5) circle (0.4);
  \node at (0,1.5) {\textbf{G}};
  \node[left=8pt] at (-0.4,1.5) {\small Générateur};
  \draw[thick] (0,1.9) -- (0,3);
  \node  at (-0.3,2) {\textbf{+}};
  \node  at (-0.3,1) {\textbf{-}};
  % === Branche 1 : lampe L1 ===
  \draw[thick] (2,3) -- (2,1.9);
  \draw[thick] (2,1.5) circle (0.4);
  \draw[thick] (1.72,1.22) -- (2.28,1.78);
  \draw[thick] (1.72,1.78) -- (2.28,1.22);
  \draw[thick] (2,1.1) -- (2,0);
  \node[right=5pt] at (2.4,1.5) {\small $\mathrm{I}_1$};
  % === Branche 2 : lampe L2 ===
  \draw[thick] (4,3) -- (4,1.9);
  \draw[thick] (4,1.5) circle (0.4);
  \draw[thick] (3.72,1.22) -- (4.28,1.78);
  \draw[thick] (3.72,1.78) -- (4.28,1.22);
  \draw[thick] (4,1.1) -- (4,0);
  \node[right=5pt] at (4.4,1.5) {\small $\mathrm{I}_2$};
  % === Nœuds ===
  \filldraw (2,3) circle (2.5pt);
  \filldraw (2,0) circle (2.5pt);
  \filldraw (4,3) circle (2.5pt);
  \filldraw (4,0) circle (2.5pt);
\end{tikzpicture}
\end{center}
}
\end{tcolorbox}

\notebox{L'ampèremètre ne perturbe pas le circuit car sa \textbf{résistance interne est très faible} (idéalement nulle). Il se comporte comme un simple fil électrique dans le circuit.}

\vfill

\textbf{6. Exemples de valeurs d'intensité courantes}

\begin{center}
\renewcommand{\arraystretch}{1.5}
\begin{tabular}{|l|c||l|c|}
\hline
\textbf{Appareil / Situation} & \textbf{Intensité} & \textbf{Appareil / Situation} & \textbf{Intensité} \\
\hline
LED & $\approx 20\,\text{mA}$ & Ampoule à incandescence (60 W) & $\approx 270\,\text{mA}$ \\
\hline
Smartphone en charge & $\approx 1\,\text{A}$ & Fer à repasser & $\approx 5\,\text{A}$ \\
\hline
Moteur d'aspirateur & $\approx 3\,\text{A}$ & Seuil de danger pour le corps humain & $\approx 10\,\text{mA}$ \\
\hline
Fusible habituel & $\approx 16\,\text{A}$ & Foudre & $\approx 30\,000\,\text{A}$ \\
\hline
\end{tabular}
\end{center}

\vfill

\textbf{7. Sécurité} \\

Un courant électrique traversant le corps humain est dangereux dès \textbf{10 mA}. Au-delà de \textbf{25 mA}, il peut provoquer une électrisation grave (tétanisation musculaire), et au-delà de \textbf{100 mA}, il peut être mortel. \\

\textbf{Remarque : } Ce n'est pas la tension (en volts) qui tue directement, mais bien l'\textbf{intensité} (en ampères) qui traverse le corps. La tension est cependant responsable de faire circuler ce courant. \\

\vfill

\newpage

\textbf{8. Exercices d'application} \\

\textbf{Exercice 1 -- Conversions}

Effectuer les conversions suivantes :
\begin{multicols}{2}
\begin{enumerate}
    \item $\mathrm{I} = 0{,}350\,\text{A}$ \quad $= \ldots\ldots\,\text{mA}$
    \item $\mathrm{I} = 45\,\text{mA}$ \quad $= \ldots\ldots\,\text{A}$
    \item $\mathrm{I} = 1{,}2\,\text{A}$ \quad $= \ldots\ldots\,\text{mA}$
    \item $\mathrm{I} = 8\,\text{mA}$ \quad $= \ldots\ldots\,\mu\text{A}$
    \item $\mathrm{I} = 500\,\mu\text{A}$ \quad $= \ldots\ldots\,\text{mA}$
    \item $\mathrm{I} = 0{,}075\,\text{A}$ \quad $= \ldots\ldots\,\text{mA}$
    \item $\mathrm{I} = 3\,200\,\text{mA}$ \quad $= \ldots\ldots\,\text{A}$
    \item $\mathrm{I} = 2{,}5\,\text{mA}$ \quad $= \ldots\ldots\,\mu\text{A}$
\end{enumerate}
\end{multicols}

\cadreblanc{25}

\textbf{Exercice 2 -- Choix du calibre}

Un ampèremètre dispose des calibres suivants : $10\,\text{A}$, $1\,\text{A}$, $200\,\text{mA}$, $20\,\text{mA}$. Quel calibre choisir pour mesurer les intensités suivantes ? Justifier.
\begin{enumerate}
    \item $\mathrm{I} = 3{,}5\,\text{A}$
    \item $\mathrm{I} = 75\,\text{mA}$
    \item $\mathrm{I} = 0{,}008\,\text{A}$
\end{enumerate}

\cadreblanc{25}

\textbf{Exercice 3 -- Circuit en série}

Dans un circuit en série composé d'un générateur, d'une lampe $L_1$ et d'une lampe $L_2$, un ampèremètre placé après $L_1$ indique $\mathrm{I}_1 = 240\,\text{mA}$.
\begin{enumerate}
    \item Quelle est l'intensité $\mathrm{I}_2$ traversant la lampe $L_2$ ?
    \item Quelle est l'intensité débitée par le générateur ?
    \item Justifier ces réponses à l'aide de la propriété du circuit série.
\end{enumerate}

\cadreblanc{25}

\textbf{Exercice 4 -- Loi des nœuds}

Dans un circuit en dérivation, le circuit principal est parcouru par une intensité $\mathrm{I} = 600\,\text{mA}$. La branche 1 est parcourue par un courant $\mathrm{I}_1 = 250\,\text{mA}$.
\begin{enumerate}
    \item Calculer l'intensité $\mathrm{I}_2$ dans la branche 2.
    \item Si on ajoute une troisième branche parcourue par $\mathrm{I}_3 = 100\,\text{mA}$, quelle serait la nouvelle intensité totale ?
\end{enumerate}

\cadreblanc{25}

\end{document}
